\begin{abstract}
\noindent
Dieser Bericht dokumentiert das Forschungsprojekt B. Im Zentrum steht die Frage, ob sich mithilfe Hessian-basierter Krümmungsmetriken vorhersagen lässt, welche Schichten eines Convolutional Neural Network (CNN) unter der Power-of-Two-Quantisierung (PoT) am stärksten geschädigt werden, wie Post-Training-Quantisierung (PTQ) und Quantization-Aware-Training (QAT) diese Krümmung dabei jeweils unterschiedlich prägen, und in welchem Ausmaß der Schaden überhaupt auf Gewichts- statt Aktivierungsquantisierung zurückgeht. Die PoT ist ein hardware-attraktives Quantisierungsschema, da sie Multiplikationen durch Bitshifts ersetzt. Gleichzeitig erzeugt sie jedoch ungleichmäßige Genauigkeitsverluste über die Schichten hinweg. Als etablierter Ansatz schlägt die HAWQ-V2-Formulierung hierfür das Produkt aus Hessian-Spur und Störungsgröße, $\mathrm{Tr}(H)\cdot\|\Delta W\|^2$, als Sensitivitäts-Score vor. Um den Schaden jeder Schicht isoliert zu bestimmen, wird jeweils nur diese eine Schicht quantisiert und der daraus resultierende Anstieg des Validierungs-Loss gemessen. Auf dieser Grundlage werden drei Scores gegen den gemessenen Schaden getestet, die rohe Krümmung, die Störungsgröße und das HAWQ-V2-Produkt, und zwar über drei Architekturen (CNN, ResNet-18, ResNet-50) sowie zwei Datensätze (CIFAR10, ImageNet100). Die Ergebnisse zeigen, dass kein Score universell zuverlässig ist. Während die rohe Krümmung die am stärksten geschädigte Schicht am konsistentesten identifiziert, wird das HAWQ-V2-Produkt durch eine Anti-Korrelation zwischen Störungsgröße und tatsächlichem Schaden systematisch in die Irre geführt. Über die reine Prädiktor-Frage hinaus behandelt der Bericht die vollständige sechsphasige Analyse-Pipeline. Diese umfasst eine Random-Init-Kontrolle, eine gescheiterte KFAC-Zerlegung, eine Dekomposition in Gewichts- und Aktivierungsanteile sowie eine Deployment-Evaluation auf realer Hardware. Dabei zeigt sich, dass die PoT-Quantisierung auf einer NVIDIA A100 keinen Laufzeitvorteil bringt, da aktuelle INT8-Kernel für Faltungsschichten fehlen, während der CPU-Pfad mit echten INT8-Kernen (fbgemm) den intendierten Hardware-Vorteil mechanistisch realisiert. Ein zentrales methodisches Ergebnis ist zudem ein Rekonziliations-Ledger, das dokumentiert, wie stark publizierte Trace-Werte von Konfigurationsentscheidungen abhängen, die in der Literatur meist unausgesprochen bleiben. Die Ergebnisse der Kernfragestellung münden schließlich in ein Workshop-Paper, das bei AXIOM @ NeurIPS 2026 eingereicht wurde.
\end{abstract}